\section{Materials and Methods}
\label{sec:methods}

\subsection{Data, candidate alignment, and count models}

We use the prepared counts, source fitness scores, and candidate identities
from the PDZ, TrpB, and NucB studies
\citep{zarin2024,johnston2024,thomas2025}. Candidate sequences are aligned
across the source scores, fitted outputs, and count features before
ranking. Source controls can contribute to a count likelihood or a
normalization total without being ranking candidates. In particular,
appended NucB controls encoding premature stops are excluded from candidate
rankings. Supplementary Section~\ref{app:reproducibility} records counts,
exclusions, and sample provenance.

Let $\data$ denote the observed counts together with their population,
replicate, timepoint, and gate identities, and let $\theta$ denote model
parameters. Conditional on the recorded total reads and a vector of
variant probabilities, each sequenced population has a multinomial
count likelihood. Assay-specific models determine these probabilities.
This accounts for the fixed-total dependence among variant counts,
but does not claim to capture every source of experimental variation.

PDZ and TrpB models transform input abundances through exponential
or sigmoid growth functions before normalizing output probabilities.
PDZ has three paired input/output biological replicates with hierarchical
growth rates. TrpB A--C uses one rate per variant across times;
D--I uses a shared input vector and hierarchical output rates.
The four-site model uses two input observations sharing a latent
composition. TrpB times are divided by the last modeled time, so
the rate scale is the assay duration rather than hours.

For NucB, each selected population's probabilities are proportional
to input abundance times the probability of exceeding its gate
threshold. The model uses survival probabilities, not probabilities
of mutually exclusive fluorescence bins. Generation-specific
populations and thresholds preserve the experimental observation
structure. Supplementary Sections~\ref{app:pdz-model}--\ref{app:nuc-model}
give the complete equations, priors, and parameterizations.

\subsection{Fitted and count-derived scores}

Write $h_i(\theta)$ for candidate $i$'s fitted parameter score and
$g_i(\data)$ for its count-derived score. In growth models without
replicate effects, $h_i=\log b_i$ for positive rate $b_i$. In the
Exponential and Lomax hierarchies, $h_i=\log m_i$ for central rate
$m_i$, a conditional mean under normal replicate variation and
a conditional median under lognormal variation. For the Gamma
inverse-rate hierarchy, $h_i=-\log\kappa_i$, with conditional mean
rate $1/\kappa_i$. For NucB, $h_i=L_i$, the fitted log-fluorescence
location. Higher scores are preferred, but the units differ
between assays.

The PDZ count score combines log enrichments relative to the unmutated
reference (wild type), using inverse-variance weights. It is a DiMSum-style approximation,
not a rerun of the source DiMSum configuration. The TrpB count
score normalizes observations using stop-codon variants and the
maximum before averaging retained timepoints and replicates.
NucB uses enrichment at the highest selected gate. The latter
two calculations reproduce the source finite scores at numerical
tolerance $10^{-12}$. Supplementary Section~\ref{app:scores}
specifies formulas, filtering, and unresolved source details.

\subsection{Fitting and sampled score collections}

We combine the count likelihood $p(\data\mid\theta)$ with a prior
$p(\theta)$, a distribution specifying parameter assumptions before
conditioning on the counts. Posterior samples target
$p(\theta\mid\data)\propto p(\data\mid\theta)p(\theta)$ using
Hamiltonian Monte Carlo (HMC). The optimized fit $\hat\theta$
also includes prior terms in its objective; despite the upstream
routine's \texttt{mle} name, it is not likelihood-only maximum
likelihood estimation. The analyses reuse stored fits from the
pinned SciStanPy implementation. Supplementary
Section~\ref{app:inference-details} distinguishes recorded job
settings, implementation defaults, and artifact-verified settings.

Let $N$ be the number of ranking-eligible candidates, $S$ the
number of retained samples in a run, and
$\boldsymbol{\fit}^{(s)}=(\fit_1^{(s)},\ldots,\fit_N^{(s)})$
the joint score vector in sample $s$. The collections are:
\begin{align}
\text{Parameter:}\quad
&\theta^{(s)}\sim p(\theta\mid\data),
&&\fit_i^{(s)}:=h_i(\theta^{(s)}),\\
\text{Optimized-fit counts:}\quad
&\data^{(s)}\sim p(\cdot\mid\hat\theta),
&&\fit_i^{(s)}:=g_i(\data^{(s)}),\\
\text{Posterior counts:}\quad
&\theta^{(s)}\sim p(\theta\mid\data),\quad
\data^{(s)}\sim p(\cdot\mid\theta^{(s)}),
&&\fit_i^{(s)}:=g_i(\data^{(s)}).
\end{align}
The count recipes condition on recorded sequencing depths and
recompute scores without refitting. They represent alternative
count outcomes under the model, not all variation in a new assay.
The score calculations also differ between parameter and count
samples, so comparisons are not an additive variance decomposition.
Table~\ref{tab:inventory} lists retained runs and samples. The
four-site TrpB landscape has optimized-fit count samples only.

\begin{table}[htbp]
\centering
\caption{Analysis inventory. A run is one landscape, model, and sampling
recipe. Retained samples are draws per run, not effective sample sizes.
The 49 posterior fits and 60 optimized fits give 158 analysis runs
because posterior fits support two sampling recipes.}
\label{tab:inventory}
\begin{tabular}{lrrr}
\toprule
Scheme & Landscapes & Runs & Retained samples \\
\midrule
Parameter & 18 & 49 & 16,000--32,000 \\
Optimized-fit counts & 19 & 60 & 967--1,000 \\
Posterior counts & 18 & 49 & 15,839--32,000 \\
\bottomrule
\end{tabular}
\end{table}

\subsection{Ranks, references, and cohorts}

Let $R_i^{(s)}$ be the descending rank of $\fit_i^{(s)}$ among
eligible candidates with finite scores in sample $s$. Rank one
is best, exact score ties receive average ranks, and an undefined
score has no rank. Write $Q_a(\fit_i)$ and $Q_a(R_i)$ for empirical
quantiles at probability $a$ over a candidate's finite samples.
The median-score reference ranks $\widetilde\fit_i:=Q_{0.50}(\fit_i)$
to obtain $\widetilde r_i$. This is distinct from the median
sampled rank $Q_{0.50}(R_i)$. Rank intervals are
$[Q_{0.05}(R_i),Q_{0.95}(R_i)]$; widths and signed endpoint
differences from $\widetilde r_i$ are stored separately.

Reported-order comparisons use the reported top 100 and top 1,000.
Rank-interval summaries use each run's median-score reference top
100. The decision cohort contains that reference batch and every
outside candidate selected into the sampled top 100 with frequency
at least 0.01. This is not a reference-top-1,000 restriction on
challengers. All displayed ranks and inclusion frequencies are
calculated against the full eligible competitor pool, not just
the displayed cohort.

We retain samples with at least 1,000 finite candidate scores.
Undefined scores are not replaced by zero. Count-based reference
medians require at least two finite scores per candidate;
parameter-based references require at least one. Eligibility can differ
slightly between score constructions because of these finite-score rules. Fixed-size
top-$k$ sets use deterministic tie selection rather than a
threshold on average ranks. Supplementary
Section~\ref{app:computation} gives the remaining finite and
tie conventions.

\subsection{Reported-rank curves and summary statistics}

Let $r_i^{\mathrm{rep}}$ be candidate $i$'s reported rank and
$\mathcal S_i$ the indices of its finite sampled ranks. Its
reported-rank sample percentile is
\[
u_i:=\frac{1}{|\mathcal S_i|}\sum_{s\in\mathcal S_i}
\left[\mathbf1\{R_i^{(s)}<r_i^{\mathrm{rep}}\}
+\tfrac12\mathbf1\{R_i^{(s)}=r_i^{\mathrm{rep}}\}\right].
\]
Here $\mathbf1\{\cdot\}$ is one when the condition holds and zero
otherwise. For each reported cohort, its empirical cumulative
distribution function is $F(u)$, and the stored area is
\[
A:=\int_0^1|F(u)-u|\,du.
\]
We integrate the step function exactly between its jumps, not
by trapezoidal interpolation through plotted vertices. Candidates
with no finite sampled rank are omitted from the curve, with
the supported cohort count recorded. Spearman correlation instead
compares $r_i^{\mathrm{rep}}$ and $Q_{0.50}(R_i)$ over the supported
cohort. Supplementary
Section~\ref{app:area-interpretation} gives contrasting cases that
have the same area and the implementation checks.

\paragraph{Mean absolute rank displacement.}
For each run, let $\mathcal C_{100}$ contain reported-top-100 candidates
with finite reported and median sampled ranks, and let
$n_{100}:=|\mathcal C_{100}|$. We calculate
\[
D_{\mathrm{rank}}:=\frac{1}{n_{100}}\sum_{i\in\mathcal C_{100}}
\left|Q_{0.50}(R_i)-r_i^{\mathrm{rep}}\right|.
\]
The median is taken over each candidate's finite sampled ranks before
the absolute difference is calculated; this is not the average absolute
displacement across individual samples. We then average candidates with
equal weight. Reported cohort membership and rank positions are taken
from the prepared tables without reranking within the cohort. All runs
have 99 or 100 supported candidates, whose counts are retained alongside
the full-precision displacement values in the exported CSV.
Table~\ref{tab:rank-displacement} aligns matching landscape/model
specifications across the parameter-posterior, optimized-fit count,
and posterior count collections, without averaging across models or
omitting specifications available under only one sampling scheme.
The table rounds to one decimal place for display. Rank ties retain
the conventions described above; the percentile equality convention
does not enter this calculation.

\subsection{Candidate inclusion and batch overlap}

Let $T_k^{(s)}$ contain the $k$ selected candidates in sample $s$.
Define $\hat q_i(k):=S^{-1}\sum_s\mathbf1\{i\in T_k^{(s)}\}$.
An undefined score counts as no selection, so the denominator
includes every retained sample. For distinct samples, the mean
shared fraction is
\[
\widehat C_k:=\frac{1}{S(S-1)}\sum_{s\ne s'}
\frac{|T_k^{(s)}\cap T_k^{(s')}|}{k}.
\]
Writing $T_k^\star$ and $T_k^{\mathrm{rep}}$ for the fixed
median-score and reported batches, we also report
\[
\widehat C_k^\star:=\frac{1}{S}\sum_s
\frac{|T_k^{(s)}\cap T_k^\star|}{k},\qquad
\widehat O_k^{\mathrm{rep}}:=\frac{1}{S}\sum_s
\frac{|T_k^{(s)}\cap T_k^{\mathrm{rep}}|}{k}.
\]
The fixed-list comparison is $|T_k^\star\cap T_k^{\mathrm{rep}}|/k$.
We use budgets $k\in\{10,20,50,100,200,500,1000\}$ and calculate
sample-pair overlap exactly from inclusion counts
(Supplementary Section~\ref{app:computation}). Two independent
uniform size-$k$ subsets of $N$ candidates share fraction $k/N$
in expectation, which provides a reference rather than a model
of engineering selection.

\subsection{Score geometry and descriptive comparisons}

Define $S_{\mathrm{run}}$ as the IQR of the top-1,000 median
scores, $W_i^{\fit}:=Q_{0.95}(\fit_i)-Q_{0.05}(\fit_i)$,
and $U_i:=W_i^{\fit}/S_{\mathrm{run}}$. The score margin is
$X_i:=(\widetilde\fit_i-\widetilde\fit_{(100)})/S_{\mathrm{run}}$,
where $\widetilde\fit_{(100)}$ is the reference rank-100 score.
The common denominator leaves within-run Spearman correlations
unchanged. Nonlinear monotone score transformations can change
widths and spacing while preserving sampled ranks, so score
geometry is interpreted on its stated scale.

Count features sum the retained observed input and output reads
with shared inputs counted once. Supplementary
Section~\ref{app:count-features} specifies the assay-dependent
sums. Within-run correlations and regressions use the decision
cohort and are descriptive associations, not estimates of the
effect of additional reads. We retain individual model-run points
and use medians across models within landscapes when summarizing
landscapes. More fitted models do not create more independent
experimental evidence.

\paragraph{Overview point displays.}
In the overview panels of Figs.~\ref{fig:reported}, \ref{fig:ranks},
\ref{fig:reported1000}, and~\ref{fig:counts}, each small filled dot
represents one scalar summary for a landscape/model/sampling-scheme
analysis, not a sample or chain. Within each landscape and scheme,
we order models by their names and run identifiers, then assign evenly
spaced vertical offsets for display. A singleton receives no additional
offset. These positions are deterministic and have no quantitative
interpretation; the horizontal metric values are not jittered. Larger
hollow symbols show medians across models at the scheme's central
vertical position. The black example outlines and arrows follow the
selected model points, including their display offsets, not the medians.

\subsection{Running examples and reproducible exports}
\label{methods:examples}

We choose the running examples to illustrate contrasting behavior,
not to represent the nearest median in each protein system. PDZ trans-1
shows many plausible replacements despite a narrow range of score widths;
TrpB libB combines low reported-rank agreement with substantial batch
turnover; NucB G3 has uncertain internal ranks but a comparatively stable
top-100 group. The two growth examples use Exponential--Exponential and
NucB uses Exponential. All three are available in each sampling collection.

We contextualize these choices using four summaries: reported-top-100
curve area and Spearman correlation, median reference-top-100 rank width,
and $\widehat C_{100}$. Within each protein system and sampling recipe,
we summarize the non-Gamma models by their median within each landscape.
None of the selected values lies outside the interval from the lower
quartile minus 1.5 IQRs to the upper quartile plus 1.5 IQRs of those
landscape summaries. This is a descriptive outlier check, not a claim
that the examples are typical on every metric. TrpB libB is near the
low-agreement end; related three-site libraries A--C also show substantial
batch instability. The complete overview and supplementary curves retain
the full context, including Gamma fits.

The example audit records all run values, their percentile positions
relative to equally weighted landscape medians, mean and largest distances
from the center, and the reason for each illustrative choice. The centrality
shortlist is descriptive and does not determine the fixed examples.
All retained posterior fits underwent sampler diagnostics; omission of a summary
from a saved output does not make a fit ineligible. The recorded numerical
diagnostics are preserved without imputing values absent from the files.

The resulting identities are passed explicitly to every main
figure and comparison table. Exports use prepared summary tables,
not a new fit or a subset of candidate points. The provenance
record contains source-table and fit-inventory hashes. Each of
the 158 analyses appears once in the point-level supplement and
once with each of the two reported cohorts in the curve supplement.
The curves share panels within each protein system and sampling scheme;
Table~\ref{tab:reported-curve-statistics} retains individual analysis
statistics. Tests check ranks, ties, cohort coverage, exact areas, example
identities, and plotted coordinates. These are implementation checks.